% !TeX encoding = UTF-8
% 第1章 绪论

\section{绪\hspace{1em}论}

知识图谱是一种语义网络知识表示形式，通过结构化方式描述现实世界中的实体及其相互关系，是人工智能领域的基础信息组织形式；在教育信息化中，如何系统组织课程知识并实现智能化应用，直接影响教学质量和学习效率，而传统知识图谱的构建依赖大量人工标注，面对专业领域的命名实体识别与关系抽取任务时，标注成本高、领域适应性差、知识更新滞后等问题突出。近年来，以GPT、DeepSeek等为代表的大语言模型（Large Language Model, LLM）在自然语言理解与生成方面进步明显，为知识抽取提供了新思路：得益于大模型在零样本和少样本场景下的学习能力，即使标注数据十分有限，也能完成较高质量的信息抽取，这为教育领域知识图谱的自动化构建创造了条件。基于上述背景，本章先综述国内外基于大语言模型的命名实体识别、关系抽取与知识图谱构建的研究现状与局限，分析课程知识抽取面临的挑战并提出研究思路，最后交代论文的主要内容和组织结构。

\subsection{研究背景与意义}

\subsubsection{研究背景与趋势}

命名实体识别（Named Entity Recognition, NER）和关系抽取（Relation Extraction, RE）是信息抽取的两个核心任务，目标是从非结构化文本中自动识别预定义类型的实体并抽取实体间的语义关系。这两项技术在知识图谱构建、智能问答、文本挖掘中应用广泛\cite{ner_survey_2024,generative_ie_survey_2023}。随着深度学习的发展\cite{zhouzhihua2016}，基于神经网络的方法在标准数据集上表现良好，但这类方法通常依赖大量高质量的标注数据进行监督学习，遇到领域迁移和新实体类型时，泛化能力不足。

教育领域中，课程知识的组织专业性强，结构紧密。以计算机网络课程为例，其中涉及大量协议实体（如TCP、UDP、IP）、概念实体（如拥塞控制、流量整形）、机制实体（如三次握手、滑动窗口）以及性能指标实体（如吞吐量、时延）。这些实体之间存在复杂的依赖和应用关系，是课程知识体系的核心。教育领域的知识抽取面临独特的挑战：专业术语存在歧义、知识表达方式多样、对证据可追溯性有要求，传统的有监督学习方法难以直接适应。

大语言模型的兴起为解决上述难题提供了新思路。近年的综述研究表明，大语言模型已从预训练语言模型发展为覆盖训练优化、能力评测和行业应用的通用技术体系\cite{wangyaozu2024_llm_survey}。通过提示工程（Prompt Engineering）和上下文学习（In-Context Learning），大语言模型可以在仅给定少量示例甚至无示例的情况下完成特定任务\cite{llms_kg_survey_2023}，这对标注数据稀缺的教育领域尤其重要。大模型较强的语义理解能力使其能捕捉文本中隐含的语义关系，涌现能力则有助于识别传统方法难以发现的隐性知识关联，而预训练阶段积累的知识表示也使其能更好地理解领域专业知识。

\subsubsection{课题目的与意义}

本课题研究利用大语言模型进行课程知识的命名实体识别与关系抽取，构建面向计算机网络课程的知识图谱自动化构建系统。具体包含以下目标：

探索适用于教育领域的大模型提示策略。通过设计任务特定的提示模板，结合零样本和少样本学习方式，提升大模型在课程知识抽取任务上的表现。针对不同的实体类型和关系类型，设计差异化的提示策略，利用大模型的上下文学习能力。通过实验对比，确定提示设计原则和参数配置。研究不同提示元素（如任务描述、示例、输出格式说明）对抽取性能的影响，为提示工程提供实践参考。

建立知识抽取的质量保障机制。引入证据验证和一致性检查机制，降低大模型生成错误知识的风险，保障知识图谱的准确性。设计自动化的质量评估指标，对抽取结果进行实时监控和反馈。建立人机协作的审核流程，兼顾效率与质量。开发冲突检测算法，识别并解决知识抽取中的逻辑矛盾和不一致问题。

设计知识图谱构建流水线。将文档解析、文本预处理、实体识别、关系抽取、知识融合等环节整合为统一流程，形成可复用的教育知识图谱构建方案。设计模块化架构，支持灵活的配置和扩展，以适应不同学科的知识抽取需求。实现从原始PDF到结构化知识图谱的端到端自动化。提供清晰的配置接口和日志监控，降低系统的使用门槛。

本课题在理论和应用两个层面均具有价值。理论上，探索大模型在专业领域知识抽取中的应用模式，为领域自适应的信息抽取方法研究提供参考。研究成果有助于理解大语言模型在特定领域的表现特征，为提示工程和领域适配方法的发展提供实证依据。通过实验设计和结果分析，探讨大模型在教育知识抽取任务中的能力边界和优化方向，为少样本和零样本场景下的信息抽取方法研究提供参考。

应用上，研究成果可直接用于课程资源的智能化管理和个性化学习推荐。构建的知识图谱可应用于智能问答、学习路径推荐、知识点关联分析等教育场景，改善教学效果和学习体验。教师可借助知识图谱快速定位课程重点，设计更有针对性的教学内容；学生可通过知识图谱建立不同知识点之间的关联，形成完整的知识体系。研究成果还可推广至其他专业领域的知识图谱构建。

\subsection{国内外研究现状}

\subsubsection{基于大模型的命名实体识别}

近年来，利用大语言模型进行命名实体识别的方法发展迅速，研究者围绕提示设计、输出约束、数据增强等方向做了大量工作。

GPT-NER\cite{gpt_ner_2023}提出了一种基于特殊标记的提示转换方法，引入\texttt{@@\#\#}等标记符号，将NER任务转化为大模型易于理解的文本生成任务，并配合自验证机制提高识别准确率。该方法解决了大模型在结构化预测任务上的格式遵循问题，使模型输出可直接解析为结构化标注。实验表明，该方法在多个标准数据集上取得了与监督学习方法相当甚至更优的成绩。研究还探索了不同标记符设计对模型性能的影响，发现特定的分隔符组合能提升实体边界识别的准确性。通过自一致性检查，降低了格式错误的发生率。

SLIMER\cite{slimer_2024}针对少样本NER场景，提出了基于实体类型定义和标注指南的提示设计框架。该方法为每种实体类型提供详细的定义描述和标注规范，引导大模型更好地理解任务要求，在跨领域NER任务上效果好。研究还发现，为实体类型添加具体示例和边界案例，能提升模型的识别准确率。通过负例说明，帮助模型理解哪些文本片段不应被识别为特定类型的实体，降低了误识别率。该方法的核心思路是利用人类标注知识指导模型学习，即使没有大量标注数据也能取得较好效果。

CodeIE\cite{codeie_2023}探索了将Python代码作为提示的新范式，利用代码的结构化特性引导大模型生成结构化的NER输出。研究表明，相比自然语言提示，代码风格的提示在结构化预测任务上具有明显优势，因为代码的语法约束能更好地引导模型生成符合格式要求的输出。该方法在零样本和少样本设置下均表现出较好的泛化能力，能够适应新的实体类型而不需要重新训练模型。代码表示还方便集成复杂的约束规则和验证逻辑，提升了抽取质量。

LLM-DA\cite{llm_da_2024}聚焦数据增强方向，提出了基于大模型的上下文级和实体级数据增强方法。通过在上下文层面进行语义变换、在实体层面进行同义词替换和边界扰动，扩充了训练数据的多样性。该方法对缓解少样本场景下的数据稀缺问题有重要参考意义，能提升模型的鲁棒性和泛化能力。研究还提出了基于大模型的数据质量评估方法，自动筛选高质量的增强样本。通过迭代增强策略逐步扩充训练数据集，在保持数据质量的同时扩大数据规模。

\subsubsection{基于大模型的关系抽取}

在关系抽取领域，大语言模型同样展现了潜力，研究者在示例选择、推理机制、任务分解等方面做了大量工作。

GPT-RE\cite{gpt_re_2023}提出了一种任务感知的示例检索策略，根据输入文本与候选示例的语义相似度动态选择上下文示例，提升了少样本关系抽取的性能。该方法通过引入实体感知的示例编码机制，使模型能更好地理解与当前任务相关的上下文信息，从而在复杂的关系分类任务中取得更好效果。实验结果表明，合理的示例选择策略对最终性能的影响甚至超过了模型规模的选择。研究还探索了动态示例库更新机制，使系统能够随着运行不断积累高质量的示例。

Revisiting-RE\cite{revisiting_re_2023}分析了链式思维（Chain-of-Thought, CoT）推理在关系抽取任务中的作用，发现显式引导大模型进行逐步推理能提高复杂关系的识别准确率，在涉及多跳推理的场景中尤为明显。研究还探讨了不同推理路径设计对最终性能的影响，可供提示工程设计参考。研究发现，关系特定的推理模板比通用模板能取得更好的效果，将复杂的关系判断分解为多个简单的推理步骤，可降低模型的判断难度。

AutoRE\cite{auto_re_2024}提出了一种基于关系-头实体-事实（Relation-Head-Fact）分解的抽取范式，将复杂的关系抽取任务分解为多个子任务，降低了模型的生成难度。该方法结合了基于人类反馈的强化学习（Reinforcement Learning from Human Feedback, RLHF）训练范式，使模型输出更契合知识图谱的结构要求。实验结果表明，分解策略能减少关系抽取中的错误传播，提升端到端的抽取质量。该方法的优势在于模块化设计，每个子任务可独立优化和评估。

\subsubsection{统一信息抽取与知识图谱构建}

统一信息抽取框架旨在通过单一模型同时完成实体识别、关系抽取等多项任务，以提升系统的效率和一致性。

InstructUIE\cite{instructuie_2023}提出了基于指令微调的多任务信息抽取框架，通过设计统一的指令格式将不同类型的信息抽取任务转化为一致的文本到文本生成问题，实现了跨任务的知识迁移。该方法在多个信息抽取任务上取得了有竞争力的性能，显示了统一框架在多任务上的优势。研究还发布了大规模的多任务指令数据集，为后续研究提供了数据基础。通过统一编码，模型能够学习不同任务之间的共享知识，提升泛化能力。

GoLLIE\cite{gollie_2023}延续了代码作为提示的思路，将标注指南编码为Python类型定义和约束条件，利用大模型的代码理解能力进行结构化信息抽取。该方法在零样本场景下表现出较好的泛化能力，能够适应新的实体类型和关系类型而不需要重新训练。代码形式的约束条件使模型能更准确地理解抽取任务的边界和要求，同时方便集成复杂的业务规则。该方法的优势在于可解释性强，便于调试和验证。

OneKE\cite{oneke_2025}提出了一种多智能体协作的知识抽取架构，通过调度多个专门化的抽取智能体分别负责实体识别、关系分类、属性抽取等子任务，再由融合模块整合各智能体的输出，提升了抽取质量。该架构的优势在于能发挥不同智能体的专长，通过协作机制取得比单一模型更好的抽取效果。系统还引入了置信度加权机制，在融合过程中优先采纳高置信度的预测结果。智能体之间的通信协议设计是该系统的主要创新。

在教育领域知识图谱构建方面，Educational-KG-Pipeline\cite{educational_kg_pipeline_2025}提出了一套面向教育资源的五阶段知识图谱构建流水线，涵盖数据采集、预处理、知识抽取、知识融合和图谱存储等环节，为教育领域的知识图谱构建提供了参考方案。Graphusion\cite{graphusion_2024}则专注于零样本知识图谱构建，提出了基于多源知识融合的图谱扩展方法，能够在无标注数据的情况下实现知识图谱的自动构建和扩展。

\subsubsection{教育领域知识图谱}

教育领域知识图谱研究近年来进展较快，研究者探索了多种构建和应用方法。

LLM-Curriculum-KG\cite{llm_curriculum_kg_2025}探索了人机协作的课程知识建模方法，将大模型的自动抽取能力与领域专家的知识校验相结合，实现了课程知识体系的构建。该方法采用人在回路（Human-in-the-Loop）模式，在保证知识准确性的同时提升了构建效率。研究还提出了面向课程知识的专门化实体类型体系和关系模式，为教育领域的知识表示提供了参考。实验结果表明，人机协作模式相比纯自动或纯人工的方式，在效率和质量间取得了平衡。

LLMs-KG-Survey\cite{llms_kg_survey_2023}对基于大语言模型的知识图谱技术进行了综述，分析了提示工程、指令微调、检索增强生成等技术在知识抽取、知识补全、知识推理等任务中的应用，梳理了当前研究的主要挑战和发展方向。综述强调领域适应性和知识质量控制在实践中的重要性，还建立了一套评估框架，用于比较不同方法在知识准确性、覆盖率和一致性等方面的表现。

\subsubsection{现有工作的局限与本课题切入点}

综上，基于大语言模型的信息抽取技术发展迅速，但在面向教育领域的课程知识抽取方面，仍有研究空白。现有工作在方法层面各有侧重：GPT-NER和SLIMER关注NER的提示设计，GPT-RE和AutoRE关注RE的任务分解，InstructUIE和GoLLIE关注统一信息抽取框架，但较少有工作将NER、RE与端到端的知识图谱构建流水线在课程知识场景下进行系统整合。在教育领域，LLM-Curriculum-KG和Educational-KG-Pipeline分别从人机协作和流水线构建的角度进行了探索，但在知识质量的可验证性、面向课程材料的专业化Schema设计以及零样本/少样本提示策略的系统对比方面仍需进一步研究。

针对上述空白，聚焦三个关键问题：一是如何设计面向课程知识的多层次实体与关系Schema，并进行零样本和少样本提示策略的系统对比；二是如何通过证据门控机制保障抽取知识的可验证性和可追溯性；三是如何构建模块化、可扩展的端到端知识图谱构建流水线。相关技术方案将在第三章和第四章中详细阐述。

\subsection{面临的挑战与研究思路}

\subsubsection{课程知识抽取面临的挑战}

尽管大语言模型在通用自然语言处理任务上表现优异，但在面向课程知识的信息抽取任务中，仍然面临以下挑战：

一是领域适应性不足。通用大语言模型知识面广，但对特定领域的专业术语和知识体系理解有限。课程知识涉及大量领域特定的概念和关系类型，需要针对教育场景进行专门的模型适配和提示设计。不同学科的知识表达方式不同，如何设计通用的领域适配框架仍是一个需要解决的问题。计算机网络这类课程专业性较强，涉及大量协议规范和技术细节，需要模型具备较强的领域理解能力。例如，同一术语在不同上下文中可能指代不同的概念，模型需要结合上下文才能准确理解。

二是输出结构化困难。大语言模型的输出本质上是自由文本，而知识图谱构建需要严格结构化的三元组数据。如何设计有效的提示策略和输出约束机制，确保模型生成的表示符合规范，是实际应用中的关键问题。在处理复杂的嵌套实体和多关系场景时，结构化的准确性直接影响后续知识融合的效果。输出中的格式错误、实体边界识别偏差、关系类型混淆等问题都会降低最终知识图谱的质量。如何设计可解析的输出格式，并处理模型偶尔产生的格式偏离，也需要在工程实现中解决。

三是知识准确性难以验证。大语言模型存在幻觉（Hallucination）现象，可能生成表面合理但实际错误的信息。在教育场景中，知识准确性尤其重要，需要建立证据验证机制，确保抽取结果的可信度和可追溯性。如何量化模型输出的置信度，在置信度偏低时触发人工审核，也是实际部署中需要权衡的问题。对于关键的专业知识点，抽取错误可能对学习者产生误导，因此需要质量控制。将抽取的知识与原文片段关联建立证据溯源，有助于提升可信度。

四是规模化处理复杂。课程材料通常以PDF、PPT等多种形式存在，包含大量图文混排内容。如何构建端到端的自动化处理流水线，将原始文档高效转换为结构化知识，是工程实现层面的主要挑战。大规模文档处理时的计算效率和成本控制也是实际应用中需要考量的因素。对于包含复杂排版、数学公式和图表说明的技术文档，准确提取其中的知识内容存在较大技术难度。如何设计可扩展的架构以支持持续的知识更新和增量处理，也是长期运营中的问题。

\subsubsection{挑战到方案的映射}

针对上述四个方面的挑战，设计了一套系统的技术方案，从领域适配、输出规整、质量保障和工程架构四个维度进行应对，如图\ref{fig:1-1}所示。

\begin{figure}[ht]
    \centering
    \begin{tikzpicture}[
        node distance=0.6cm and 0.5cm,
        challenge/.style={
            rectangle, rounded corners=2pt, draw=black!70, dashed, fill=black!4,
            text width=2.8cm, minimum height=1cm, align=flush center,
            font=\footnotesize, text=black!85, inner xsep=4pt
        },
        solution/.style={
            rectangle, rounded corners=2pt, draw=black!70, fill=black!8,
            text width=2.8cm, minimum height=1cm, align=flush center,
            font=\footnotesize, text=black!85, inner xsep=4pt
        },
        center/.style={
            rectangle, rounded corners=4pt, draw=black, very thick, fill=black!15,
            text width=2.4cm, minimum height=1.2cm, align=center,
            font=\small\bfseries, text=black!90, inner xsep=4pt
        },
        arr_l/.style={-{stealth}, thick, draw=black!70},
        arr_r/.style={-{stealth}, thick, draw=black!70}
    ]
    % Left: Challenges
    \node[challenge] (c1) {领域适应性不足\\{\scriptsize 专业术语理解有限，领域特征表达多样}};
    \node[challenge, below=0.2cm of c1] (c2) {输出结构化困难\\{\scriptsize 自由文本与结构化三元组的需求矛盾}};
    \node[challenge, below=0.2cm of c2] (c3) {知识准确性难以验证\\{\scriptsize 大模型幻觉问题，教育场景可信度要求高}};
    \node[challenge, below=0.2cm of c3] (c4) {规模化处理复杂\\{\scriptsize PDF/PPT等多源文档的端到端自动化挑战}};

    % Center
    \node[center, right=1.0cm of c2] (sys) {课程知识\\[0.1cm]抽取系统};

    % Right: Solutions
    \node[solution, right=1.0cm of sys] (s1) {领域Schema与提示策略（第三章）\\{\scriptsize 四类实体与关系体系，零/少样本设计}};
    \node[solution, below=0.2cm of s1] (s2) {两阶段抽取架构（第三章）\\{\scriptsize Propose-Review流程，JSON结构化约束}};
    \node[solution, below=0.2cm of s2] (s3) {证据门控机制（第三章）\\{\scriptsize 原文子串匹配验证，可追溯的证据溯源}};
    \node[solution, below=0.2cm of s3] (s4) {六阶段流水线架构（第四章）\\{\scriptsize 模块化设计，断点续跑，独立扩展}};

    % Arrows: challenges → center (fan pattern, all converge on sys.west)
    \draw[arr_l] (c1.east) -- (sys.west);
    \draw[arr_l] (c2.east) -- (sys.west);
    \draw[arr_l] (c3.east) -- (sys.west);
    \draw[arr_l] (c4.east) -- (sys.west);

    % Arrows: center → solutions (fan pattern, all diverge from sys.east)
    \draw[arr_r] (sys.east) -- (s1.west);
    \draw[arr_r] (sys.east) -- (s2.west);
    \draw[arr_r] (sys.east) -- (s3.west);
    \draw[arr_r] (sys.east) -- (s4.west);

    % Labels (consistent 0.2cm spacing)
    \node[below=0.2cm of c4, font=\footnotesize, text=black!60] (clabel) {面临挑战};
    \node[below=0.2cm of s4, font=\footnotesize, text=black!60] (slabel) {解决思路};
    \end{tikzpicture}
    \caption{课程知识抽取面临的挑战与解决思路}\label{fig:1-1}
\end{figure}

具体而言：针对领域适应性不足，定义了覆盖协议、概念、机制和性能指标的四类实体类型体系，以及包含、依赖、对比和应用于的四类关系类型体系，结合零样本和少样本两种提示策略提升LLM的领域理解能力；针对输出结构化困难，设计了Propose-Review两阶段抽取架构，通过审核阶段的类型合法性检查、证据子串匹配和实体存在性验证等多项约束，将LLM的自由文本输出规整为结构化三元组；针对知识准确性难以验证，提出了证据门控机制，要求每条三元组必须附带原始文本块的连续原文子串作为依据，实现了输出结果的自动验证与可追溯；针对规模化处理复杂，构建了文档渲染、OCR文字提取、文本分块、知识抽取、知识融合和图数据库导入的六阶段端到端流水线，各阶段以文件系统中间产物解耦，支持独立扩展与断点续跑。这些技术方案的详细设计将在第三章和第四章中展开。

\subsection{论文的主要内容与结构}

\subsubsection{论文主要内容}

本课题围绕基于大语言模型的课程知识命名实体识别与关系抽取展开研究，构建了面向计算机网络课程的知识图谱自动化构建系统。主要研究内容包括：

\textbf{（1）多源课程语料库构建}

构建了包含205份PDF文档的多源课程语料库，涵盖教材、课件、实验指导、RFC文档等多种类型。设计了标准化的文档预处理流程，实现从PDF到结构化文本的自动转换，最终生成16,438个文本块，总计约1,040万词元（tokens），为后续的知识抽取提供了数据基础。针对不同来源的文档特性，设计了差异化的解析策略，保障文本提取的完整性和准确性。对于包含复杂排版的技术文档，采用了专门的处理方法以保留知识结构信息。建立了文档元数据管理系统，记录每份文档的来源、类型和处理状态。

\textbf{（2）基于大模型的命名实体识别}

针对课程知识的特性，定义了协议、概念、机制、性能指标四类实体类型。设计了零样本和少样本两种提示策略，通过任务特定的提示模板引导大模型识别文本中的命名实体。实验表明，该方法在没有大规模标注数据的情况下依然能达到较高的识别准确率。针对不同实体类型的特点，设计了差异化的提示模板和示例选择策略，提升了识别效果。引入了实体链接机制，将识别出的实体与标准化术语库关联，解决了实体歧义和同义词问题。

\textbf{（3）基于大模型的关系抽取}

定义了包含、依赖、对比、应用于四类关系类型。提出了“提议-审核”两阶段抽取策略，第一阶段由大模型生成候选关系三元组，第二阶段通过证据验证机制确保抽取结果的准确性和可追溯性。引入了关系约束规则，对不符合逻辑的关系进行过滤，降低了错误率。设计了关系置信度评估机制，为后续的知识融合提供质量依据。对于复杂的多跳关系，设计了迭代抽取策略，逐步构建关系链。建立了关系类型层次结构，支持细粒度的关系分类。

\textbf{（4）知识抽取流水线设计}

设计了涵盖文档解析、图像转换、文本识别、文本分块、知识抽取、知识融合、图谱导入七个环节的自动化处理流水线。各模块之间通过标准化接口进行数据交换，支持灵活的配置和扩展。实现了基于配置文件的流水线编排，用户可根据需求自定义处理流程和参数设置。流水线支持断点续处理和增量更新，提高了大规模数据处理的效率。提供了详细的日志记录和进度监控功能，便于跟踪处理状态。设计了错误处理机制，对处理失败的文档自动记录并重试。

\textbf{（5）应用验证与系统实现}

在计算机网络课程数据上进行了应用验证，最终构建的知识图谱（\texttt{fs-recheck}配置下的\texttt{merged-fs-recheck}）包含32,760个实体和17,627个关系三元组。其中，高频实体包括TCP协议（2,117次）、UDP协议（1,069次）、HTTP协议（701次）、IP协议（640次）等，准确反映了课程知识的核心内容。基于500条独立人工金标在严格、软匹配、槽位与类型级四档口径下完成评估，验证了系统在课程知识抽取任务上的可用性。通过多个典型应用场景的测试，验证了知识图谱在教学辅助和学习推荐中的应用潜力。

\subsubsection{论文结构}

上述研究内容按照“技术基础、系统设计、系统实现、实验评估”的逻辑展开，从相关技术分析出发，经方案设计与代码实现，最终通过实验验证系统的有效性。论文共分为六章，各章的主要内容如下：

第一章为绪论，阐述课题的研究背景与意义，综述国内外研究现状与局限，分析课程知识抽取面临的挑战并提出研究思路，最后介绍论文的主要内容与结构安排。

第二章为相关技术基础，从大语言模型基础出发，依次介绍命名实体识别技术、关系抽取技术、知识图谱构建与实体规范化技术，以及教育领域知识图谱的相关工作，为后续章节的系统设计与实现提供技术支撑。

第三章为课程知识抽取系统设计，从需求分析入手，依次阐述实体与关系Schema设计、文档预处理设计概要、文本分块设计、实体识别与关系抽取的架构设计思想、知识融合策略以及图存储方案，重点说明各模块的设计决策与取舍理由。

第四章为系统实现，首先给出六阶段流水线的总体架构，随后详细阐述文档渲染与OCR、文本分块、命名实体识别、关系抽取、知识融合以及图存储等各模块的工程实现细节，包括关键参数配置、算法实现与代码架构。

第五章为实验与结果分析，报告实验环境配置和课程语料构建结果，从实体抽取、关系抽取、零样本与少样本对比、知识图谱质量分析以及典型案例等多个角度对系统进行实验评估。

第六章为总结与展望，总结课题的主要研究工作与贡献，分析当前工作的局限，并对未来的研究方向进行展望。
